\begin{abstract}
    This work proposes a controlled numerical study of sparse linear
    systems arising from a semi-implicit finite-difference
    discretization of the two-dimensional linearized shallow water
    equations. Rather than developing a full oceanographic simulation,
    the equations are used as a model problem to investigate how
    spatial resolution and variations in water depth affect the
    conditioning of the discrete system and the convergence of
    iterative solvers.

    The experiments will consider a constant-depth configuration
    and spatially heterogeneous depth profiles with controlled
    coefficient contrast. The Conjugate Gradient method without
    preconditioning will be compared with Jacobi-preconditioned
    and incomplete-Cholesky-preconditioned variants, provided that
    the resulting matrix is verified to be symmetric positive
    definite. The evaluation will consider iteration counts,
    relative residual norms, estimated condition numbers, and
    computational time. Preconditioner setup costs and the potential
    reuse of preconditioners across successive time steps will also
    be investigated.

    The methodological contribution is a controlled, cost-aware
    comparison of these solver strategies across different grid
    resolutions and coefficient configurations. The study aims to
    identify the conditions under which preconditioning improves
    overall computational performance, rather than assessing its
    effectiveness solely through iteration counts.
\end{abstract}

\noindent\textbf{Keywords:} Shallow water equations;
semi-implicit discretization; sparse linear systems;
iterative solvers; preconditioning.